\clearpage\fancyhead[L]{\small KOMET.}\hypertarget{native-350-lower}{}\label{native-350-lower}
\begin{center}\textbf{KOMET.}\par\vspace{4pt}\hrule\vspace{4pt}\fontsize{10}{12}\selectfont Göttingische gelehrte Anzeigen. Heft 15, S. 147..149, 25. Januar 1812.\par\vspace{4pt}\hrule\vspace{4pt}\end{center}

\par\indent Seit den ersten in diesen Blättern (Heft~201 des Vorjahres [S.~342 dieses Bandes]) angekündigten Beobachtungen des \emph{neuen} Kometen hat das seit mehr als einem Monat ungewöhnlich ungünstige Wetter nur drei weitere Beobachtungen an der hiesigen Sternwarte zugelassen. Wir teilen sie zusammen mit den ersten Beobachtungen mit, deren Reduktion etwas verbessert wurde.

\begin{center}
\fontsize{10}{12}\selectfont\setlength{\tabcolsep}{3pt}\renewcommand{\arraystretch}{1.15}\begin{tabular}{l|c|c|c}
 & \shortstack{Mittlere\\Zeit} & \shortstack{Rektaszension} & Deklination \\
\hline
1811 Dez. 9 & $10^{\text{u}}6^{\text{m}}52^{\text{s}}$ & $63^\circ49'41''4$ & $10^\circ21'55''5$ S. \\
11 & $10\;\;34\;\;1$ & $63\;\;33\;\;18.0$ & $8\;\;39\;\;46.4$ \\
12 & $8\;\;5\;\;52$ & $63\;\;26\;\;25.3$ & $7\;\;54\;\;25.9$ \\
1812 Januar 3 & $7\;\;1\;\;30$ & $62\;\;37\;\;6.4$ & $9\;\;47\;\;14.7$ N. \\
4 & $6\;\;53\;\;59$ & $62\;\;40\;\;38.6$ & $10\;\;29\;\;57.8$ \\
\end{tabular}
\end{center}

\noindent Schon am 13. Dezember hatte Herr Professor \textsc{Gauß} aus seinen ersten Beobachtungen und denen des Herrn von \textsc{Zach}, die wir ebenfalls dort mitgeteilt haben, vorläufige parabolische Elemente berechnet. Sie sind im Westphälischen Moniteur, Nr.~301, und im Dezemberheft der Monatlichen Correspondenz abgedruckt. Natürlich konnten diese auf so dürftige Daten gestützten Elemente nur Näherungen sein, doch betrug ihre Abweichung von der Beobachtung am 3. Januar noch immer nur 8 Minuten. Nachdem Herr Professor \textsc{Gauß} seine Beobachtung vom 4. Januar gemacht hatte, übertrug er Herrn \textsc{Nicolai} die Verbesserung der parabolischen Elemente. Dieser führte die Arbeit mit größter Sorgfalt und mit bestem Erfolg aus. Die\SourceLeaf{351} folgenden verbesserten Elemente hat dieser geschickte Rechner ermittelt. Sie werden sich gewiss nur noch geringfügig verbessern lassen.

\Needspace{15\baselineskip}
\begin{center}\fontsize{10}{12}\selectfont\setlength{\tabcolsep}{4pt}\renewcommand{\arraystretch}{1.15}\begin{tabular}{l}
Zeit des Periheldurchgangs 1811 Nov. 11\\
\quad $9^{\mathrm u}\;24^{\mathrm m}\;38^{\mathrm s}$ Meridian von Göttingen.\\
\end{tabular}\end{center}
\begin{center}\fontsize{10}{12}\selectfont\setlength{\tabcolsep}{3pt}\renewcommand{\arraystretch}{1.15}\begin{tabular}{llr}
Abstand im Perihel & .\;.\;.\;.\;.\;.\;.\;.\;.\;.\;. & $1.589198$\\
Länge des Perihels & .\;.\;.\;.\;.\;.\;.\;.\;.\;.\;.\;.\;. & $47^\circ\;39'\;36''1$\\
Länge des aufsteigenden Knotens & .\;.\;.\;.\;.\;.\;.\;.\;. & $92\;53\;44.2$\\
Neigung der Bahn & .\;.\;.\;.\;.\;.\;.\;.\;.\;.\;.\;.\;.\;. & $31\;32\;38.7$\\
\end{tabular}\end{center}
\noindent Rechtläufige Bewegung.

\par\indent Die Konstanten, mit denen nach der Methode von Herrn Professor \textsc{Gauß} die Koordinaten des Kometen in Bezug auf den Äquator berechnet werden, sind die folgenden.

\Needspace{13\baselineskip}
\[
\begin{aligned}
x\cos\tfrac{1}{2}\upsilon^2&=\alpha\sin(\upsilon+138^\circ\;9'\;39''6)\\
y\cos\tfrac{1}{2}\upsilon^2&=\ObservedBeta\sin(\upsilon+59\;53\;53.5)\\
z\cos\tfrac{1}{2}\upsilon^2&=\gamma\sin(\upsilon+355\;26\;5.4)\\
\log\alpha&=0.1319477\\
\log\ObservedBeta&=0.1784770\\
\log\gamma&=9.9865984
\end{aligned}
\]

\Needspace{35\baselineskip}\par\indent Außer den oben angeführten Beobachtungen erhielt Herr Professor \textsc{Gauß} noch einige Beobachtungen von Herrn von \textsc{Lindenau} an der Seeberger Sternwarte, von Herrn Dr. \textsc{Olbers} in Bremen und von Herrn \textsc{Burckhardt} in Paris. Burckhardts Beobachtungen sind Meridianbeobachtungen. Mit all diesen Beobachtungen verglich Herr \textsc{Nicolai} seine Elemente und erhielt die folgende schöne Übereinstimmung, die kaum noch etwas zu wünschen übrig lässt.

\begin{center}
\fontsize{10}{12}\selectfont\setlength{\tabcolsep}{5pt}\renewcommand{\arraystretch}{1.15}\begin{tabular}{l|c|c|l}
 & \multicolumn{2}{c|}{Unterschied} & \\
 & \shortstack{in Rektaszension} & \shortstack{in Deklination} & Beobachter \\
\hline
1811 Nov. 18 & $-12''4$ & $-138''4$ & v. Zach \\
19 & $+34.7$ & $+19.5$ & v. Zach \\
20 & $-11.9$ & $-5.5$ & v. Zach \\
21 & $-45.0$ & $+4.9$ & v. Zach \\
Dez. 8 & $+9.7$ & .\;.\;.\;. & v. Lindenau \\
9 & $+1.0$ & $-51.4$ & Olbers \\
9 & $+4.9$ & $+22.6$ & v. Lindenau \\
9 & $-5.0$ & $+3.4$ & Gauss \\
11 & $+2.9$ & $-32.7$ & Gauss \\
12 & $+8.5$ & $-27.6$ & Gauss \\
14 & $+23.5$ & $-34.4$ & Olbers \\
14 & $+18.9$ & $-6.7$ & Burckhardt \\
16 & $+28.7$ & $-18.1$ & Olbers \\
22 & $+20.8$ & $+2.4$ & Burckhardt \\
25 & $+25.1$ & $-3.1$ & Burckhardt \\
1812 Jan. 3 & $+2.7$ & $-21.1$ & Olbers \\
3 & $-5.5$ & $+12.9$ & Gauss \\
4 & $+2.8$ & $-12.7$ & Gauss \\
6 & $-35.6$ & $-1.5$ & Olbers \\
\end{tabular}
\end{center}

\par\indent Seit den ersten hiesigen Beobachtungen, bei denen seine Helligkeit am größten und ungefähr $= 0.7$ war, ist der Komet merklich lichtschwächer geworden. Doch ließ er sich am 3. und 4. Januar noch ebenso gut wie zuvor beobachten. Er würde noch sehr lange sichtbar bleiben, wenn nicht zu befürchten\SourceLeaf{352} wäre, dass ihn sein immer schwächer werdendes Licht wohl schon im Februar unseren Blicken entziehen wird. Inzwischen werden die Beobachtungen durch die folgende Ephemeride sehr erleichtert, die ebenfalls Herr \textsc{Nicolai} berechnet hat. Ihr zufolge wird der Komet, nachdem er bereits mitten durch die Hyaden gezogen ist, aus dem Stier in den Fuhrmann treten und dort wahrscheinlich unsichtbar werden.
